\documentclass[11pt]{article}
\usepackage[margin=1in]{geometry}
\usepackage{booktabs}
\usepackage{hyperref}

\title{Sealed-Bid Sourcing}
\author{AthenaTheOwl}
\date{June 2026}

\begin{document}
\maketitle

\begin{abstract}
This draft defines a small benchmark for reverse sourcing with sealed bids.
The v0.1 artifact contains a typed 10 supplier, 5 lot scenario, a reference
sealed runtime, an unsealed comparison runtime, and a surplus report. The
sealed runtime publishes the scenario hash, supplier commitments, winners,
opened winning bids, and per-lot scores. It does not publish losing bid
values. The construction is scoped to private set intersection, additive
secret sharing, and hash commitments as cited in the bibliography
\cite{sciencedirect2026psi,freedman2004private,decristofaro2010practical,shamir1979share,beaver1991efficient}.
\end{abstract}

\section{Introduction}

Reverse auctions rely on suppliers believing that a submitted bid is used for
award evaluation and not for later price extraction. Buyer-side comparison
agents weaken that belief when losing bid prices are disclosed to the buyer.
This repository gives the claim a runnable target: compare a sealed reference
runtime with an unsealed baseline on the same procurement scenario.

\section{Threat Model}

The v0.1 sealed runtime assumes an honest-but-curious buyer and a
non-colluding audit verifier. It excludes collusion between those parties,
malicious denial of service, and production side-channel resistance. The
runtime leaks only the fields named in \texttt{paper/leakage\_bounds.yaml}.

\section{Scenario}

The canonical scenario is \texttt{scenarios/v0\_10x5.json}. It contains ten
suppliers, five lots, per-lot scoring weights, reserve prices, capacity
constraints, supplier defensive markup parameters for the unsealed baseline,
and bid attributes for price, lead time, ESG score, and quality score.

\section{Reference Runtime}

The sealed runtime evaluates committed bid payloads and opens the winning bid
for each lot. The checked-in implementation is a reference implementation. It
models the audit surface and receipt shape; it is not production MPC code.

\section{Unsealed Baseline}

The unsealed runtime applies each supplier's defensive markup before scoring.
This models the procurement setting where suppliers expect buyer-side agents to
observe losing prices and therefore stop submitting their best price.

\section{Results}

The checked-in report is \texttt{reports/surplus\_delta.md}. It compares buyer,
supplier, and total surplus between the two receipts emitted from the canonical
scenario.

\section{Limitations}

The benchmark is intentionally small. It does not claim throughput, production
cryptographic hardening, native e-sourcing integration, or a malicious-buyer
security proof.

\bibliographystyle{plain}
\bibliography{refs}

\end{document}
